% ============================================================================
% Section 5: Agent Mutation Testing
% ============================================================================

\section{Agent Mutation Testing}
\label{sec:mutation}

Mutation testing assesses test suite quality by measuring its ability to
detect injected faults. For traditional software, mutants are syntactic
source code modifications. For AI agents, the ``source'' comprises
prompts, tool configurations, models, and context---requiring entirely
new mutation operators. This section defines four classes of
agent-specific mutation operators, formalizes the mutation score under
stochastic semantics, and proves a mutation adequacy theorem.

\subsection{Agent Mutation Operators}
\label{sec:mutation:operators}

\begin{definition}[Agent Mutation Operator]
\label{def:mutation-operator}
An \emph{agent mutation operator} is a function $m: \mathcal{A} \to
\mathcal{A}$ that transforms an agent configuration into a \emph{mutant}
agent $A' = m(A)$ by modifying exactly one component of
$A = (\pi, \mathcal{T}, \mu, \omega)$.
\end{definition}

We define four classes of mutation operators, each targeting a different
component.

\subsubsection{Prompt Mutation Operators ($M_{\text{prompt}}$)}
\label{sec:mutation:prompt}

\begin{definition}[Prompt Mutations]
\label{def:prompt-mutations}
The class $M_{\text{prompt}}$ consists of four operators:
\begin{enumerate}[leftmargin=*]
  \item \textbf{Synonym substitution} ($m_{\text{syn}}$): Replace a key
    instruction word with a synonym. For example,
    ``\texttt{Summarize the document}'' $\to$
    ``\texttt{Condense the document}''.

  \item \textbf{Instruction reordering} ($m_{\text{reorder}}$): Permute
    the order of instructions within the prompt, preserving content but
    altering priority. If $\pi = (i_1, i_2, \ldots, i_k)$, then
    $m_{\text{reorder}}(\pi) = (i_{\sigma(1)}, \ldots,
    i_{\sigma(k)})$ for a random permutation~$\sigma$.

  \item \textbf{Noise injection} ($m_{\text{noise}}$): Insert irrelevant
    or distracting text into the prompt. Models the effect of prompt
    pollution or context contamination.

  \item \textbf{Instruction dropout} ($m_{\text{drop}}$): Remove a
    randomly selected instruction from the prompt. Models the effect of
    incomplete or truncated prompts.
\end{enumerate}
\end{definition}

\begin{remark}
Prompt mutations target the \emph{semantic resilience} of the agent: a
well-designed agent should be robust to minor prompt perturbations
(synonym substitution, reordering) but sensitive to meaningful changes
(instruction dropout). The mutation score thus measures the test suite's
ability to distinguish semantic from non-semantic changes.
\end{remark}

\subsubsection{Tool Mutation Operators ($M_{\text{tool}}$)}
\label{sec:mutation:tool}

\begin{definition}[Tool Mutations]
\label{def:tool-mutations}
The class $M_{\text{tool}}$ consists of three operators:
\begin{enumerate}[leftmargin=*]
  \item \textbf{Tool removal} ($m_{\text{remove}}$): Remove a randomly
    selected tool from $\mathcal{T}$, producing
    $\mathcal{T}' = \mathcal{T} \setminus \{t_j\}$. Tests whether the
    agent gracefully degrades or fails silently.

  \item \textbf{Tool reordering} ($m_{\text{treorder}}$): Permute the
    order of tools in the tool inventory (relevant when tool descriptions
    are presented to the LLM as a list). Models the position bias
    exhibited by some language models.

  \item \textbf{Tool noise} ($m_{\text{tnoise}}$): Modify a tool's
    description or schema by injecting misleading information. Models
    the effect of tool description errors or adversarial tool
    definitions~\citep{clawhavoc2026}.
\end{enumerate}
\end{definition}

\subsubsection{Model Mutation Operators ($M_{\text{model}}$)}
\label{sec:mutation:model}

\begin{definition}[Model Mutations]
\label{def:model-mutations}
The class $M_{\text{model}}$ consists of two operators:
\begin{enumerate}[leftmargin=*]
  \item \textbf{Model swap} ($m_{\text{swap}}$): Replace the underlying
    LLM $\mu$ with a different model from the same capability tier.
    For example, replace GPT-4o with Claude 3.5 Sonnet. Tests
    cross-model compatibility.

  \item \textbf{Version downgrade} ($m_{\text{version}}$): Replace the
    model with an older or less capable version (e.g., GPT-4o $\to$
    GPT-4o-mini). Tests graceful degradation under model changes.
\end{enumerate}
\end{definition}

\subsubsection{Context Mutation Operators ($M_{\text{context}}$)}
\label{sec:mutation:context}

\begin{definition}[Context Mutations]
\label{def:context-mutations}
The class $M_{\text{context}}$ consists of three operators:
\begin{enumerate}[leftmargin=*]
  \item \textbf{Context truncation} ($m_{\text{trunc}}$): Truncate the
    input context to $\lfloor \gamma \cdot |x| \rfloor$ tokens for
    $\gamma \in (0, 1)$. Models context window overflow.

  \item \textbf{Context noise} ($m_{\text{cnoise}}$): Inject irrelevant
    documents or passages into the context. Models the effect of noisy
    retrieval in RAG pipelines.

  \item \textbf{Context permutation} ($m_{\text{cperm}}$): Permute the
    order of context segments. Models sensitivity to document ordering
    in retrieval-augmented settings.
\end{enumerate}
\end{definition}

\subsection{Stochastic Kill Semantics}
\label{sec:mutation:kill}

In traditional mutation testing, a mutant is ``killed'' if the test
produces a different output. For stochastic agents, outputs differ across
runs even without mutations. We must therefore define kill semantics that
distinguish \emph{mutation-induced} behavioral changes from
\emph{natural stochastic variation}.

\begin{definition}[Stochastic Mutation Kill]
\label{def:kill}
Let $A$ be the original agent and $A' = m(A)$ a mutant. Let
$\mathcal{S}$ be a test suite with scenarios $\{S_1, \ldots, S_k\}$.
The mutant $A'$ is \emph{killed} by $\mathcal{S}$ if \emph{at least
one} of the following conditions holds for some scenario $S_j$:
\begin{enumerate}[leftmargin=*]
  \item \textbf{Verdict change:} $V(A, S_j) \neq V(A', S_j)$ where $V$
    is the stochastic verdict function (\cref{def:verdict}).

  \item \textbf{Significant score difference:} The pass rates differ
    significantly:
    \begin{equation}
    \label{eq:kill-score}
      |\hat{p}(A, S_j) - \hat{p}(A', S_j)| \geq \delta_{\min}
      \quad\text{and}\quad
      p\text{-value}(H_0: p_A = p_{A'}) < \alpha_{\text{kill}}
    \end{equation}
    where $\delta_{\min}$ is a minimum effect size threshold and
    $\alpha_{\text{kill}}$ is the kill significance level.

  \item \textbf{Distributional shift:} The distribution of outputs
    (or trace features) differs significantly, as measured by a
    two-sample test (e.g., Kolmogorov-Smirnov for continuous metrics
    or $\chi^2$ for categorical outcomes).
\end{enumerate}
\end{definition}

\begin{remark}
Condition (1) is the strongest: the test suite's verdict flips from
\PASS{} to \FAIL{} (or vice versa). Condition (2) is more sensitive: it
detects changes even when the verdict remains the same. Condition (3) is
the most general: it catches distributional shifts in agent behavior
(e.g., the agent now takes longer paths or uses different tools) even
if the pass rate is unchanged.
\end{remark}

\begin{definition}[Equivalent Mutant]
\label{def:equivalent-mutant}
A mutant $A' = m(A)$ is \emph{equivalent} if, for all inputs
$x \in \mathcal{X}$ and all sample sizes $n$, the output distributions
of $A$ and $A'$ are identical: $A(x) \overset{d}{=} A'(x)$. Equivalent
mutants cannot be killed by any test suite and are excluded from the
mutation score.
\end{definition}

The equivalent mutant problem is undecidable in general. For agents, we
use a practical heuristic: a mutant is \emph{presumed equivalent} if,
after $n_{\text{equiv}}$ trials on each of $k_{\text{equiv}}$ scenarios,
no statistically significant difference is detected at level
$\alpha_{\text{equiv}}$.

\subsection{Mutation Score}
\label{sec:mutation:score}

\begin{definition}[Agent Mutation Score]
\label{def:mutation-score}
Let $\mathcal{M}(A) = \{A'_1, \ldots, A'_q\}$ be the set of mutants
generated from agent $A$ and $\mathcal{M}_{\text{equiv}} \subseteq
\mathcal{M}(A)$ the set of equivalent mutants. The \emph{agent mutation
score} of test suite $\mathcal{S}$ is:
\begin{equation}
\label{eq:mutation-score}
  \text{MS}(\mathcal{S}, A) =
    \frac{|\{A'_i \in \mathcal{M}(A) \setminus \mathcal{M}_{\text{equiv}}
      : A'_i \text{ is killed by } \mathcal{S}\}|}
         {|\mathcal{M}(A) \setminus \mathcal{M}_{\text{equiv}}|}
\end{equation}
\end{definition}

\begin{remark}[Class-Specific Mutation Scores]
We also report mutation scores per operator class:
$\text{MS}_{\text{prompt}}, \text{MS}_{\text{tool}},
\text{MS}_{\text{model}}, \text{MS}_{\text{context}}$. These
fine-grained scores identify which aspects of the agent are well-tested
and which require additional test scenarios.
\end{remark}

\subsection{Mutation Adequacy}
\label{sec:mutation:adequacy}

\begin{theorem}[Mutation Adequacy]
\label{thm:mutation-adequacy}
Let $\mathcal{S}$ be a test suite with mutation score
$\text{MS}(\mathcal{S}, A) \geq \tau$ for threshold $\tau \in (0, 1]$.
Let $\delta > 0$ be a behavioral impact threshold. If the mutation
operators satisfy the \emph{coupling effect} assumption---that complex
faults are coupled to simple faults detectable by the mutation
operators---then for any real regression $A \to A^*$ with behavioral
impact $\Delta(A, A^*) \geq \delta$:
\begin{equation}
\label{eq:adequacy-bound}
  \Pr[\mathcal{S} \text{ detects the regression}] \geq
    \tau \cdot \left(1 - e^{-\delta / \delta_0}\right)
\end{equation}
where $\delta_0$ is a characteristic scale parameter of the mutation
operator suite, determined by the average behavioral impact of a single
mutation.
\end{theorem}

\begin{proof}[Proof sketch]
The proof proceeds by contrapositive. Suppose the test suite fails to
detect a regression of magnitude $\delta$. By the coupling effect, this
regression can be decomposed into a composition of simple mutations.
Since the test suite achieves mutation score $\tau$, it detects a
fraction $\tau$ of all simple mutations. The probability that \emph{none}
of the constituent simple mutations are detected decreases exponentially
with $\delta / \delta_0$, yielding the bound. The full proof, including
the formal coupling assumption and its justification for agent systems,
is in \cref{app:proof:mutation}.
\end{proof}

\begin{corollary}
\label{cor:mutation-complete}
If $\text{MS}(\mathcal{S}, A) = 1$ (all non-equivalent mutants killed)
and the mutation operators are \emph{complete} (every behavioral change
of magnitude $\geq \delta$ includes at least one mutation from
$\mathcal{M}$), then $\Pr[\mathcal{S} \text{ detects regression}] \to 1$
as $\delta / \delta_0 \to \infty$.
\end{corollary}

\subsection{Cost-Aware Mutation Testing}
\label{sec:mutation:cost}

Running a full mutation analysis is expensive: for $q$ mutants, each
tested over $n$ trials on $k$ scenarios, the total cost is
$O(q \cdot n \cdot k)$ agent invocations. We introduce two cost
reduction strategies.

\begin{definition}[Selective Mutation]
\label{def:selective-mutation}
Rather than generating all possible mutants, \emph{selective mutation}
generates mutants only from a curated subset of operators identified as
most effective. Let $\mathcal{M}^* \subset \mathcal{M}$ be the selected
operator subset. The \emph{selective mutation score} is:
\begin{equation}
\label{eq:selective-ms}
  \text{MS}^*(\mathcal{S}, A) =
    \frac{|\text{killed}(\mathcal{S}, \mathcal{M}^*(A))|}
         {|\mathcal{M}^*(A) \setminus \mathcal{M}_{\text{equiv}}|}
\end{equation}
\end{definition}

\begin{definition}[SPRT-Accelerated Mutation Kill]
\label{def:sprt-mutation}
Rather than running $n$ fixed trials to determine if a mutant is killed,
apply the SPRT (\cref{def:sprt}) to stop as soon as the kill/survive
decision is clear. Formally, test $H_0: p_{A'} = p_A$ vs.\
$H_1: |p_{A'} - p_A| \geq \delta_{\min}$ sequentially, terminating when
the log-likelihood ratio crosses a boundary.
\end{definition}

\begin{proposition}[Mutation Cost Reduction]
\label{prop:mutation-cost}
Combining selective mutation (selecting a
$\gamma = |\mathcal{M}^*|/|\mathcal{M}|$ fraction of operators) with
SPRT\nobreakdash-accelerated kills (achieving expected savings
factor~$\sigma$) reduces total mutation testing cost by:
\begin{equation}
\label{eq:cost-reduction}
  \text{Cost reduction} = 1 - \gamma \cdot \sigma
\end{equation}
With $\gamma = 0.3$ (30\% of operators) and $\sigma = 0.5$ (50\% fewer
trials via SPRT), the total cost is reduced by 85\%.
\end{proposition}

\subsection{Mutation Operator Design Principles}
\label{sec:mutation:principles}

We identify four principles for designing effective agent mutation
operators:

\begin{enumerate}[leftmargin=*]
  \item \textbf{Semantic granularity.} Mutations should produce
    behavioral changes at varying granularity levels---from subtle
    (synonym substitution) to severe (model swap)---mirroring the
    spectrum of real-world agent changes.

  \item \textbf{Component isolation.} Each mutation should modify exactly
    one component, enabling precise attribution of detected regressions
    to specific agent aspects.

  \item \textbf{Ecological validity.} Mutations should model realistic
    changes that occur in practice: prompt edits, tool updates, model
    upgrades, context changes. Arbitrary random perturbations are less
    informative.

  \item \textbf{Composability.} Individual mutations should compose to
    model complex changes. If mutations $m_1$ and $m_2$ are individually
    meaningful, then $m_2 \circ m_1$ should model a realistic compound
    change.
\end{enumerate}
